% ===================================================================
%  CUADRO N.º 12 — La Estación. — Los Ferrocarriles. — Los Barcos.
%  Traduction 2026 d'apres texto/io, controlee sur texto/fr et
%  texto/en. Consigne : voir texto/es/10-cuadro-01.tex.
% ===================================================================

%%K t12-apar-1 apar
\VUsaut{4.0mm}
\VUtitre{15.0pt}{60}{CUADRO N.º 12}
\VUsaut{2.4mm}
\VUpk{11.4pt}{40}{La Estación. --- Los Ferrocarriles. --- Los Barcos.}
\VUsaut{3.4mm}
\textsc{Nota.} --- \textit{Los horarios usados en cada país son distintos,
por lo que hemos tomado como ejemplo las horas del} \VUgras{cuadro
francés}.

%%K t12-01-1 p
1. --- Acabo de hacer un viaje por Alemania. Antes de emprenderlo había
comprado una \VUgras{guía de horarios} \textsuperscript{(15)} para saber la hora de
salida de los trenes. Habiendo comprobado que el tren más cómodo salía de
París a las nueve de la noche y llegaba a Estrasburgo a la una y trece,
preparé mi viaje y, al día siguiente, me hice llevar a la estación.

%%K t12-02-1 p
2. --- Cuando entré en la gran sala de salidas, detrás de un viejo
\VUgras{cura} \textsuperscript{(2)} de pueblo que llevaba en la mano su \VUgras{bolsa
de viaje} \textsuperscript{(3)}, ya habían llegado muchos \VUgras{viajeros}
\textsuperscript{(1)}.

%%K t12-02-2 p
Como quería mandar un \VUgras{telegrama} \textsuperscript{(5)}, hice que un
\VUgras{revisor} \textsuperscript{(6)} me indicara la \VUgras{oficina de telégrafos}
\textsuperscript{(4)}.

%%K t12-03-1 p
3. --- Antes de pasar a la \VUgras{sala de espera} \textsuperscript{(7)} fui a sacar
un \VUgras{billete} \textsuperscript{(9)} de ida y vuelta.

%%K t12-04-1 p
4. --- No esperé mucho en la \VUgras{taquilla} \textsuperscript{(8)}, porque había
poca gente. Un \VUgras{soldado} \textsuperscript{(10)} que salía de permiso me cedió
amablemente su turno, de modo que tuve tiempo de pedir algunos datos al
\VUgras{jefe de estación} \textsuperscript{(13)}, que hablaba con el \VUgras{comisario}
\textsuperscript{(12)} de vigilancia y con el \VUgras{gendarme} \textsuperscript{(11)} de
servicio.

%%K t12-tit-1 sub
\VUpk{10.2pt}{40}{Facturar el equipaje.}
\VUsaut{3.0mm}

%%K t12-05-1 p
5. --- Después de comprar un periódico en el \VUgras{quiosco de libros}
\textsuperscript{(14)}, hice pesar mi \VUgras{equipaje} \textsuperscript{(17)}, que un
\VUgras{mozo} \textsuperscript{(16)} acababa de traer en una \VUgras{carretilla}
\textsuperscript{(19)}; el \VUgras{empleado} \textsuperscript{(22)} encargado de la
\VUgras{báscula} \textsuperscript{(21)} me dijo que mi baúl pesaba treinta y cinco
kilos, lo que hacía un \VUgras{exceso} de cinco kilos, por el que tenía que
pagar un suplemento en la \VUgras{ventanilla de equipajes} \textsuperscript{(23)}.

%%K t12-06-1 p
6. --- Como iba con adelanto, tuve ocasión de mirar el movimiento de los
viajeros por la estación. Unos dejaban o recogían pequeños bultos en la
\VUgras{consigna} \textsuperscript{(25)}; otros consultaban los \VUgras{horarios}
\textsuperscript{(26)}. Los viajeros que llegaban se apresuraban hacia la
\VUgras{salida} \textsuperscript{(32)} con su \VUgras{maleta} \textsuperscript{(24)} en la
mano. Algunos esperaban en la \VUgras{entrega de equipajes} \textsuperscript{(27)} y
presentaban su \VUgras{resguardo} \textsuperscript{(29)} para retirar sus baúles,
\VUgras{cajas} \textsuperscript{(28)} o \VUgras{cestas} \textsuperscript{(30)}.

%%K t12-07-1 p
7. --- Los \VUgras{empleados de consumos} \textsuperscript{(33)} registraban con
cuidado los bultos para ver si había algo que declarar.

%%K t12-08-1 p
8. --- Algunos viajeros se dirigían a la \VUgras{fonda} \textsuperscript{(34)}, donde
un \VUgras{camarero} \textsuperscript{(35)} se apresuraba a servirlos. Tenían que
darse prisa, porque el \VUgras{tren} \textsuperscript{(36)}, que era un expreso, no se
quedaba mucho tiempo en la estación.

%%K t12-09-1 p
9. --- Pasé luego al andén de salida y busqué un \VUgras{vagón}
\textsuperscript{(37)} directo para no tener que cambiar durante el viaje. Subí a
un coche de pasillo que tiene un \VUgras{departamento} \textsuperscript{(38)} para
fumadores y otro reservado a «señoras solas».

%%K t12-tit-2 sub
\VUpk{10.2pt}{40}{Salida de noche.}
\VUsaut{3.0mm}

%%K t12-10-1 p
10. --- Después de subir el \VUgras{estribo} \textsuperscript{(40)}, cerrar la
\VUgras{portezuela} \textsuperscript{(39)}, subir la \VUgras{cortinilla} \textsuperscript{(41)} y
colocar mis bultos pequeños en la \VUgras{rejilla} \textsuperscript{(43)}, me senté en
el \VUgras{asiento} \textsuperscript{(42)}. Al ver al \VUgras{farolero} \textsuperscript{(44)}
encender las lámparas, pensé que tenía una noche que pasar en el vagón y
llamé al empleado encargado del alquiler de las \VUgras{almohadas}
\textsuperscript{(45)} para pedirle una.

%%K t12-11-1 p
11. --- El conductor del tren vino a comprobar los billetes. Le pregunté
por qué no salíamos, puesto que era la hora. Me contestó que el
\VUgras{jefe de tren} \textsuperscript{(46)} estaba ocupado haciendo colocar unas
bicicletas en el \VUgras{furgón} \textsuperscript{(47)}. Además, todavía no se habían
cambiado todos los \VUgras{calientapiés} \textsuperscript{(48)}.

%%K t12-12-1 p
12. --- Pero no íbamos a tardar en salir, porque el \VUgras{fogonero}
\textsuperscript{(50)}, sobre su \VUgras{ténder} \textsuperscript{(49)}, cogía carbón para
avivar el fuego de la \VUgras{locomotora} \textsuperscript{(54)}, y \VUgras{el
maquinista} \textsuperscript{(51)} solo esperaba la señal del \VUgras{jefe adjunto de
estación} \textsuperscript{(52)}, que estaba en el \VUgras{andén} \textsuperscript{(53)}.

%%K t12-13-1 p
13. --- La \VUgras{caldera} \textsuperscript{(56)} de la locomotora rugía sordamente;
la \VUgras{chimenea} \textsuperscript{(55)} vomitaba torrentes de humo mezclado con
\VUgras{vapor} \textsuperscript{(60)}. Sonó un \VUgras{silbido} \textsuperscript{(59)} agudo y
los \VUgras{pistones} \textsuperscript{(58)} se pusieron en movimiento, accionando las
\VUgras{ruedas} \textsuperscript{(57)} de la pesada máquina.

%%K t12-tit-3 sub
\VUsaut{3.0mm}
\VUpk{10.2pt}{40}{¡En marcha!}
\VUsaut{3.0mm}

%%K t12-14-1 p
14. --- La velocidad del tren, que corría sobre los \VUgras{raíles}
\textsuperscript{(64)}, se aceleró; los \VUgras{andenes} \textsuperscript{(65)} se acabaron;
pasamos los \VUgras{retretes} \textsuperscript{(66)} y también una fila de vagones
apartados en otra \VUgras{vía} \textsuperscript{(63)}: \VUgras{coches cama}
\textsuperscript{(67)}, \VUgras{coche restaurante} \textsuperscript{(68)}, \VUgras{coche correo}
\textsuperscript{(69)}.

%%K t12-noto-1 noto
\VUnotes{143.2mm}{%
\textit{(*) Nota.} --- \textit{Se entiende que el viaje se describe tal
como era a lo largo de la frontera de antes de la guerra.}%
}

%%K t12-15-1 p
15. --- Luego pasó ante nuestros ojos la \VUgras{estación de mercancías}
\textsuperscript{(71)}. Bajo los cobertizos, unos empleados cargaban las
\VUgras{mercancías} \textsuperscript{(72)} enviadas en pequeña velocidad. Unos
\VUgras{mozos de maniobras} \textsuperscript{(76)}, junto al \VUgras{depósito de agua}
\textsuperscript{(73)}, maniobraban \VUgras{vagones de ganado} \textsuperscript{(75)} y
\VUgras{vagones plataforma} \textsuperscript{(79)} para formar un \VUgras{tren de
mercancías} \textsuperscript{(80)}.

%%K t12-16-1 p
16. --- Pasamos la última \VUgras{aguja} \textsuperscript{(78)}, junto a la cual
estaba el guardagujas; dejamos atrás el \VUgras{disco de señales}
\textsuperscript{(86)} y la estación se perdió en la lejanía.

%%K t12-17-1 p
17. --- Pronto entramos en el \VUgras{puente de hierro} \textsuperscript{(74)} que
cruza el \VUgras{río} \textsuperscript{(81)}. Pasado el puente, seguimos el curso del
río, por el que potentes \VUgras{remolcadores} \textsuperscript{(82)} arrastraban
pesadas \VUgras{gabarras} \textsuperscript{(83)}. Vimos además un \VUgras{barco de
pasajeros} \textsuperscript{(84)} en el momento en que iba a atracar en el
\VUgras{pontón} \textsuperscript{(85)}.

%%K t12-18-1 p
18. --- Después de pasar de largo varias estaciones, atravesamos algunos
\VUgras{túneles} \textsuperscript{(87)} y dejamos tras nosotros innumerables
\VUgras{casetas} \textsuperscript{(88)} de \VUgras{guardabarreras}. Mecido por el
traqueteo del tren, me quedé profundamente dormido.

%%K t12-tit-4 sub
\VUsaut{3.0mm}
\VUpk{10.2pt}{40}{La estación de Deutsch-Avricourt.}
\VUsaut{3.0mm}

%%K t12-19-1 p
19. --- Me despertó de pronto la voz de un empleado que gritaba:
«¡Deutsch-Avricourt! \textsuperscript{(*)} ¡Bajen todos!». Era la inspección de
aduanas y todas las fastidiosas formalidades de la frontera. Tuve que
poner mi reloj en hora con el \VUgras{reloj} \textsuperscript{(89)} de la estación,
que iba cincuenta y cinco minutos adelantado; apenas despierto,
distinguía con dificultad la \VUgras{aguja grande} \textsuperscript{(90)} de la
\VUgras{pequeña} \textsuperscript{(91)}; la \VUgras{esfera} \textsuperscript{(92)} misma estaba
del todo confusa por la niebla de la mañana.

%%K t12-20-1 p
20. --- Mientras los aduaneros hacían su inspección, descifré bostezando
los \VUgras{carteles} \textsuperscript{(93)} que decoran las paredes de la estación.
Desayuné para hacer tiempo, consulté el mapa de la \VUgras{red}
\textsuperscript{(94)} alemana para ver qué distancia me separaba todavía de
Estrasburgo, y volví a subir al vagón, reconfortado por la esperanza de
llegar pronto al término de mi viaje.
\VUsaut{6.0mm}
\VUfilet{22mm}
